\chapter{Carbon nanotubes}
\label{ch:cnt}

\paragraph{In this chapter.} Carbon allotropes, SWNT vs MWNT, the three fabrication methods compared, the CVD growth
mechanism and chirality.

\section{Structure}
Carbon allotropes \slide{105}: diamond, graphite, fullerene (C$_{60}$), nanotube, graphene.
A nanotube is a \textbf{rolled-up graphene sheet}: one sheet gives a \textbf{single-walled} nanotube (SWNT); rolling
graphite (several sheets) gives a \textbf{multi-walled} nanotube (MWNT), wall spacing $\approx\SI{0.34}{nm}$ \slide{106, 115}.

\section{Fabrication methods compared}
\begin{center}
\renewcommand{\arraystretch}{1.25}
\begin{tabularx}{\linewidth}{@{}lXXX@{}}
\toprule
 & \textbf{Arc vaporisation} \slide{107} & \textbf{Laser vaporisation} \slide{108} & \textbf{CVD} \slide{109--119} \\ \midrule
Principle & Arc between two C electrodes $\sim$1\,mm apart in inert gas; C deposits on cathode & Pulsed laser on a graphite
target in a furnace at \SI{1200}{\celsius}, Ar flow; tubes collected on a water-cooled Cu collector &
Hydrocarbon gas (C$_x$H$_y$) + inert gas over a \textbf{metal catalyst} (Ni, Co, Fe) at \SIrange{500}{1000}{\celsius} \\
Products & SWNT + MWNT (+ other products) & \textbf{High-purity SWNT} & Long SWNT/MWNT, aligned forests \\
Size control & None & Narrow distribution & Set by catalyst particle size \\
Yield & $\eta\approx30\,\%$ & Low production yield & $\eta\approx97\,\%$ \\
Other & Simplest; length $\sim\SI{50}{\micro m}$ & Expensive & Catalyst eases reaction, selects products, can be recovered \\
\bottomrule
\end{tabularx}
\end{center}
\begin{classbox}
CVD: the inert gas is there ``to not oxide [oxidise] hydrocarbon gases''. The catalyst ``eases the reaction, selects
products, [can be] recovered'' \slide{109}.
\end{classbox}

\section{CVD growth mechanism}
\slide{113--117}:
\begin{enumerate}
  \item The hydrocarbon decomposes on the metal particle; C dissolves (Fe $\to$ FeC$_x$ carbide) and diffuses.
  \item C precipitates as graphitic walls.
  \item \textbf{Tip growth} if the metal--support interaction is weak (particle lifted at the tip);
        \textbf{base growth} if it is strong (particle stays on the support).
\end{enumerate}
\begin{keybox}[: catalyst size sets tube type \slide{115}]
Catalyst particle \SIrange{1}{5}{nm} $\to$ SWNT; \SIrange{20}{50}{nm} $\to$ MWNT; $>\SI{50}{nm}$ $\to$ nano-onions
(encapsulated particle). ``Depending on nanoparticle sizes we can modify nanotube sizes.''
Catalyst film thickness/concentration sets the \emph{density} of tubes: thick film $\to$ large, dense catalyst
islands; thin film $\to$ small, sparse ones \slide{117}.
\end{keybox}
Example \slide{110}: Fe/Al$_2$O$_3$ (1 g) in C$_2$H$_6$/H$_2$ at \SI{680}{\celsius} for 2 h gives a carbon yield of
$\sim6000$ wt\,\% of the Fe weight. Carbon nanofibres improve the mechanical properties of graphite cloth \slide{112}.
Fullerenes can be inserted in tubes (``peapods'') and fuse under the electron beam \slide{121}.

\section{Chirality}
\slide{120}: the roll-up vector $(n,m)$ of the graphene sheet defines \textbf{zig-zag} $(n,0)$, \textbf{armchair}
$(n,n)$ and \textbf{chiral} $(n,m)$ tubes. The structure determines the \textbf{electronic, optical and mechanical}
properties (e.g.\ metallic vs semiconducting tubes).
